\documentclass[11pt,a4paper]{article}
\usepackage[utf8]{inputenc}
\usepackage[T1]{fontenc}
\usepackage[spanish,provide=*]{babel}
\addto\captionsspanish{\renewcommand{\tablename}{Tabla}}
\usepackage{lmodern}
\usepackage{microtype}
\usepackage[margin=2.5cm]{geometry}
\usepackage{amsmath,amssymb}
\usepackage{graphicx}
\usepackage{booktabs}
\usepackage{adjustbox}
\usepackage{pdflscape}
\usepackage{float}
\usepackage[font=small,labelfont=bf]{caption}
\usepackage{enumitem}
\usepackage{xcolor}
\usepackage{listings}
\usepackage[round]{natbib}
\usepackage[hidelinks]{hyperref}
\setlength{\parskip}{0.5em}
\setlength{\parindent}{0pt}
\setlist{itemsep=0.2em}

% Texto en rojo = algo que hay que revisar/actualizar después de volver a correr el do-file.
% Para sacar el color: \renewcommand{\revisar}[1]{#1}
\newcommand{\revisar}[1]{\textcolor{red}{#1}}
% Inserta una tabla exportada desde Stata (se achica si no entra)
\newcommand{\tabla}[1]{\begin{adjustbox}{max width=\linewidth, max totalheight=0.72\textheight}\input{#1}\end{adjustbox}}
% Figura que todavía no existe (se genera al correr el do-file): muestra un recuadro
\newcommand{\figura}[2][0.75\linewidth]{\IfFileExists{#2}{\includegraphics[width=#1]{#2}}{\fbox{\parbox{\dimexpr#1-2\fboxsep-2\fboxrule\relax}{\centering\vspace{2em}\textcolor{red}{Figura pendiente: se genera al correr el do-file}\\[0.5em]\texttt{\detokenize{#2}}\vspace{2em}}}}}
% Notas al pie de tablas
\newcommand{\notas}[1]{\par\vspace{0.3em}\begin{minipage}{0.92\linewidth}\footnotesize\textit{Notas:} #1\end{minipage}}
% Consignas en cursiva
\newcommand{\consigna}[1]{\par\textit{#1}\par}

\lstset{
  language={},
  basicstyle=\ttfamily\scriptsize,
  breaklines=true,
  columns=fullflexible,
  keepspaces=true,
  showstringspaces=false,
  frame=single,
  rulecolor=\color{gray!50},
  commentstyle=\color{gray},
  morecomment=[l]{*\ },
  morecomment=[l]{//},
  morecomment=[s]{/*}{*/},
  tabsize=4,
  inputencoding=utf8,
  extendedchars=true,
  literate={á}{{\'a}}1 {é}{{\'e}}1 {í}{{\'i}}1 {ó}{{\'o}}1 {ú}{{\'u}}1
           {Á}{{\'A}}1 {É}{{\'E}}1 {Í}{{\'I}}1 {Ó}{{\'O}}1 {Ú}{{\'U}}1
           {ñ}{{\~n}}1 {Ñ}{{\~N}}1 {ü}{{\"u}}1 {Ü}{{\"U}}1
           {¿}{{?`}}1 {¡}{{!`}}1 {’}{{'}}1 {‘}{{`}}1 {“}{{``}}1 {”}{{''}}1
           {ã}{{\~a}}1 {õ}{{\~o}}1 {â}{{\^a}}1 {ê}{{\^e}}1 {ô}{{\^o}}1 {à}{{\`a}}1
           {Ã}{{\~A}}1 {Õ}{{\~O}}1 {Â}{{\^A}}1 {Ê}{{\^E}}1 {Ô}{{\^O}}1 {Ç}{{\c C}}1 {×}{{$\times$}}1
           {°}{{$^\circ$}}1 {–}{{--}}1 {—}{{---}}1 {·}{{$\cdot$}}1 {ç}{{\c c}}1
}
\begin{document}

\begin{center}
{\Large\bfseries Problem Set 9: Power calculations}\\[0.5em]
E178: Economía Aplicada --- Universidad de San Andrés\\
Fecha de entrega: 15 de noviembre de 2024\\[0.5em]
\textbf{Integrantes:} Juan Batalla, Nina Di Costanzo Pereira, Catalina Gallardo e Ignacio Yenarópulos
\end{center}

Siguiendo la clase tutorial, simulamos una población de 15,000 empresas con ganancias $G_i\sim N(10000,2000^2)$ e impuestos pagados $Y_i=0.2\,G_i+\varepsilon_i$. Para cada tamaño de muestra (de 1,000 a 10,000) y cada efecto (1\%, 2.5\%, 5\%, 7.5\% y 10\% de aumento en los impuestos de los tratados), sacamos 500 muestras, asignamos el tratamiento al azar, estimamos $Y_i=\alpha+\tau T_i+u_i$ con errores robustos y calculamos el poder como la proporción de muestras en las que rechazamos $H_0:\tau=0$ con un test a una cola al 5\% ($t>1.65$).

Los resultados se pueden entender con la fórmula del poder. Si $\sigma$ es el desvío del término de error, $N$ el tamaño de muestra y $p$ la proporción de tratados, el error estándar del estimador es
\[
SE(\hat\tau)=\frac{\sigma}{\sqrt{N\,p\,(1-p)}}\qquad\text{y}\qquad \text{Poder}\approx\Phi\!\left(\frac{\tau}{SE(\hat\tau)}-z_{0.95}\right),
\]
con $z_{0.95}=1.65$. El poder aumenta con el tamaño del efecto y con $N$, disminuye con $\sigma$, y es máximo cuando $p=0.5$.

\section*{1. Gráfico de la tutorial}

\begin{figure}[H]\centering
\includegraphics[width=0.72\linewidth]{"../output/Figura 1.png"}
\caption{Poder según el tamaño de muestra y el tamaño del efecto (desvío del error = 500).}\label{fig:1}
\end{figure}

La Figura~\ref{fig:1} muestra que el poder aumenta con el número de observaciones, porque el error estándar cae a una tasa $1/\sqrt{N}$, y que es mayor cuanto más grande es el efecto, porque un efecto mayor es más fácil de distinguir del ruido. Con efectos de 5\% o más el poder es prácticamente 1 a partir de 2,000--3,000 observaciones. En cambio, para detectar un efecto de 2.5\% con un poder de 80\% se necesitan alrededor de 4,000 observaciones, y con un efecto de 1\% el poder no llega a 50\% ni siquiera con 10,000 observaciones. Los valores simulados coinciden con la fórmula. Por ejemplo, con un efecto de 2.5\% (50 pesos sobre impuestos promedio de 2,000) y 2,000 observaciones, el desvío de $Y$ es $\sqrt{400^2+500^2}\approx640$, el error estándar es $640/\sqrt{2000\times0.25}\approx28.6$ y el poder teórico es $\Phi(50/28.6-1.65)\approx0.54$, igual al de la simulación.

\section*{2. Varianza del error igual a 5,000}

\begin{figure}[H]\centering
\includegraphics[width=0.72\linewidth]{"../output/Figura 2.png"}
\caption{Poder con desvío del error igual a 5,000.}\label{fig:2}
\end{figure}

Al aumentar la dispersión del término de error, el outcome tiene mucho más ruido que no explica el tratamiento, el error estándar del estimador aumenta y el poder cae para cualquier tamaño de muestra y de efecto (Figura~\ref{fig:2}). Ahora un efecto de 5\% necesita alrededor de 6,000 observaciones para alcanzar un poder de 80\%, y los efectos de 1\% y 2.5\% tienen poder muy bajo incluso con 10,000 observaciones.

Hay dos detalles del código (que mantuvimos igual al de la tutorial) que conviene aclarar. Primero, \texttt{rnormal(0,5000)} fija el \emph{desvío} estándar del error en 5,000, no la varianza. Segundo, al eliminar las observaciones con impuestos negativos se descarta alrededor de un tercio de la población, y las que quedan tienen impuestos promedio de unos 4,800, con un desvío de unos 3,400. Como el efecto es proporcional a los impuestos, el efecto en pesos es mayor que en la simulación original, y por eso el poder no cae tanto como sugeriría multiplicar el desvío del error por diez. Con estos valores la fórmula vuelve a reproducir la simulación: un efecto de 10\% con 1,000 observaciones tiene un poder teórico de 0.72 (0.68 en la simulación).

\section*{3. Asignación del tratamiento al 20\% y al 80\% de las observaciones}

\begin{figure}[H]\centering
\begin{minipage}{0.49\linewidth}\centering
\includegraphics[width=\linewidth]{"../output/Figura 3.png"}\\
(a) 20\% de tratados
\end{minipage}\hfill
\begin{minipage}{0.49\linewidth}\centering
\includegraphics[width=\linewidth]{"../output/Figura 4.png"}\\
(b) 80\% de tratados
\end{minipage}
\caption{Poder con desvío del error igual a 5,000 y distintas proporciones de tratados.}\label{fig:34}
\end{figure}

Con el mismo tamaño de muestra, asignar el tratamiento a una proporción distinta de 50\% reduce el poder (Figura~\ref{fig:34}, comparada con la Figura~\ref{fig:2}). El motivo es que la varianza del estimador es proporcional a $1/[p(1-p)]$: con $p=0.5$, $p(1-p)=0.25$, mientras que con $p=0.2$ o $p=0.8$, $p(1-p)=0.16$, de modo que el error estándar es un 25\% mayor. Intuitivamente, la diferencia de medias se estima con menos precisión porque uno de los dos grupos es chico. Por ejemplo, con un efecto de 5\% y 2,000 observaciones el poder cae de alrededor de 0.48 a 0.35. Como $p(1-p)$ es simétrico, asignar el 20\% o el 80\% al tratamiento da teóricamente el mismo poder, y las dos figuras son, en efecto, muy parecidas; las pequeñas diferencias se deben al azar de la simulación (y a que, con un efecto proporcional, la varianza del grupo tratado es apenas mayor).

\section*{4. Agregar las ganancias estimadas como control}

\begin{figure}[H]\centering
\includegraphics[width=0.72\linewidth]{"../output/Figura 5.png"}
\caption{Poder con desvío del error igual a 500, controlando por las ganancias estimadas.}\label{fig:5}
\end{figure}

Al controlar por las ganancias estimadas el poder aumenta respecto de la Figura~\ref{fig:1} (Figura~\ref{fig:5}). Las ganancias explican una parte importante de la variación de los impuestos ($0.2\,G_i$ tiene un desvío de 400), así que al incluirlas como control esa variación deja de estar en el término de error. El desvío residual cae de $\approx640$ a 500, el error estándar de $\hat\tau$ se reduce alrededor de un 22\% y el estadístico $t$ aumenta en la misma proporción. Como el tratamiento es aleatorio, el control no está correlacionado con $T$, de modo que no cambia el valor esperado de $\hat\tau$, solo lo hace más preciso. Por ejemplo, con un efecto de 2.5\% y 2,000 observaciones el poder sube de 0.53 a alrededor de 0.68, y con un efecto de 1\% y 10,000 observaciones de 0.48 a 0.62.

\newpage
\appendix
\section*{Apéndice: código}
\lstinputlisting{../programs/ps9.do}

\end{document}
